% glossary.tex -- Module 7 ESSENTIALS glossary (A-Z): the key terms only.
% src/mod7/glossary.tex (term, chapter number, meaning and Chinese copied word for
% word). Edit the script, not this file.
\unnumberedchapter[\faBook]{Glossary}

The key words, A to Z, each with the chapter that explains it and the usual Chinese term.

{\small
\begin{glossarytable}
  % ---- A
  \gterm{Absorption spectrum\ {\footnotesize\textcolor{hscmuted}{\mdseries(1)}}}{a continuous spectrum crossed by dark lines; made when light passes through a cooler, low-density gas}{吸收光谱}
  \gterm{Aether (luminiferous aether)\ {\footnotesize\textcolor{hscmuted}{\mdseries(1, 4)}}}{an invisible material once thought to fill space and carry light; it does not exist}{（光）以太}
  \gterm{Annihilation\ {\footnotesize\textcolor{hscmuted}{\mdseries(4)}}}{a particle and its antiparticle meet and turn completely into photons}{湮灭}
  \gterm{Antiparticle\ {\footnotesize\textcolor{hscmuted}{\mdseries(4)}}}{a particle with the same mass as another particle but the opposite charge}{反粒子}
  % ---- B
  \gterm{Black body\ {\footnotesize\textcolor{hscmuted}{\mdseries(1, 3)}}}{an ideal object that absorbs all radiation that hits it and is a perfect emitter; its spectrum depends only on its temperature}{黑体}
  \gterm{Blueshift\ {\footnotesize\textcolor{hscmuted}{\mdseries(1)}}}{spectral lines moved to shorter wavelengths, because the source moves towards us}{蓝移}
  % ---- C
  \gterm{Coherent\ {\footnotesize\textcolor{hscmuted}{\mdseries(2)}}}{waves with the same frequency and a fixed phase difference (they stay in step)}{相干的}
  \gterm{Constructive interference\ {\footnotesize\textcolor{hscmuted}{\mdseries(2)}}}{waves meet in phase and add to a bigger wave (bright light)}{相长干涉}
  \gterm{Continuous spectrum\ {\footnotesize\textcolor{hscmuted}{\mdseries(1)}}}{a spectrum with all wavelengths and no gaps, from a hot, dense object}{连续光谱（连续谱）}
  \gterm{Corpuscle\ {\footnotesize\textcolor{hscmuted}{\mdseries(2)}}}{Newton's name for a tiny particle of light}{微粒（光微粒）}
  % ---- D
  \gterm{Destructive interference\ {\footnotesize\textcolor{hscmuted}{\mdseries(2)}}}{waves meet \ang{180} out of phase and cancel (darkness)}{相消干涉}
  \gterm{Diffraction\ {\footnotesize\textcolor{hscmuted}{\mdseries(2)}}}{the spreading of a wave as it passes through a gap or around an edge}{衍射}
  \gterm{Diffraction grating\ {\footnotesize\textcolor{hscmuted}{\mdseries(1, 2)}}}{a plate with thousands of equally spaced, parallel slits; it spreads light into a spectrum}{衍射光栅}
  \gterm{Doppler effect\ {\footnotesize\textcolor{hscmuted}{\mdseries(1)}}}{the change in wavelength when a source and an observer move towards or away from each other}{多普勒效应}
  % ---- E
  \gterm{Electromagnetic (EM) spectrum\ {\footnotesize\textcolor{hscmuted}{\mdseries(1)}}}{all EM waves, in order of frequency or wavelength}{电磁波谱}
  \gterm{Electromagnetic (EM) wave\ {\footnotesize\textcolor{hscmuted}{\mdseries(1)}}}{a travelling wave of electric and magnetic fields that change together}{电磁波}
  \gterm{Electron volt (eV)\ {\footnotesize\textcolor{hscmuted}{\mdseries(3)}}}{the energy an electron gains moving through \qty{1}{V}: \qty{1}{\eV} $=$ \qty{1.602e-19}{J}}{电子伏特}
  \gterm{Emission spectrum\ {\footnotesize\textcolor{hscmuted}{\mdseries(1)}}}{bright lines on a dark background, from an excited, low-density gas}{发射光谱（明线光谱）}
  % ---- F
  \gterm{Frame of reference\ {\footnotesize\textcolor{hscmuted}{\mdseries(4)}}}{a viewpoint, with its own rulers and clocks, for measuring positions, times and motion}{参考系}
  % ---- H
  \gterm{Huygens' principle\ {\footnotesize\textcolor{hscmuted}{\mdseries(2)}}}{every point on a wavefront sends out wavelets; the line touching them is the new wavefront}{惠更斯原理}
  % ---- I
  \gterm{Inertial frame of reference\ {\footnotesize\textcolor{hscmuted}{\mdseries(4)}}}{a frame that is not accelerating: at rest or moving at constant velocity}{惯性参考系}
  \gterm{Intensity ($I$)\ {\footnotesize\textcolor{hscmuted}{\mdseries(2, 3)}}}{the power of light falling on each square metre (\si{W.m^{-2}}); for one colour, it is set by the number of photons each second}{光强（强度）}
  \gterm{Interference\ {\footnotesize\textcolor{hscmuted}{\mdseries(2)}}}{the pattern made when overlapping waves add or cancel}{干涉}
  % ---- L
  \gterm{Length contraction\ {\footnotesize\textcolor{hscmuted}{\mdseries(4)}}}{a moving object is measured to be shorter along its direction of motion}{长度收缩（尺缩效应）}
  \gterm{Line broadening\ {\footnotesize\textcolor{hscmuted}{\mdseries(1)}}}{spectral lines becoming wider, from fast rotation or a dense atmosphere}{谱线展宽}
  % ---- M
  \gterm{Mass defect ($\Delta m$)\ {\footnotesize\textcolor{hscmuted}{\mdseries(4)}}}{mass before minus mass after a process: the mass that becomes energy}{质量亏损}
  \gterm{Muon\ {\footnotesize\textcolor{hscmuted}{\mdseries(4)}}}{an unstable particle like a heavy electron, made by cosmic rays high in the atmosphere}{μ子（缪子）}
  % ---- P
  \gterm{Path difference ($\delta$)\ {\footnotesize\textcolor{hscmuted}{\mdseries(2)}}}{how much further one wave travels than the other}{波程差（路程差）}
  \gterm{Photoelectric effect\ {\footnotesize\textcolor{hscmuted}{\mdseries(3)}}}{electrons are released from a metal surface when light shines on it}{光电效应}
  \gterm{Photon\ {\footnotesize\textcolor{hscmuted}{\mdseries(1, 3)}}}{one packet (quantum) of light energy, $E = hf$}{光子}
  \gterm{Plane-polarised light\ {\footnotesize\textcolor{hscmuted}{\mdseries(2)}}}{light whose electric field vibrates in one direction only}{线偏振光（平面偏振光）}
  \gterm{Polariser\ {\footnotesize\textcolor{hscmuted}{\mdseries(2)}}}{a filter that passes only the vibration along its transmission axis}{起偏器（偏振片）}
  \gterm{Positron ($\mathrm{e}^+$)\ {\footnotesize\textcolor{hscmuted}{\mdseries(4)}}}{the antiparticle of the electron: same mass, positive charge}{正电子}
  \gterm{Postulate\ {\footnotesize\textcolor{hscmuted}{\mdseries(4)}}}{a starting statement that a theory takes as true}{基本假设（公设）}
  \gterm{Proper length ($l_0$)\ {\footnotesize\textcolor{hscmuted}{\mdseries(4)}}}{the length of an object measured in the frame where it is at rest}{固有长度（原长）}
  \gterm{Proper time ($t_0$)\ {\footnotesize\textcolor{hscmuted}{\mdseries(4)}}}{the time between two events, measured by a clock that is present at both events}{固有时（原时）}
  % ---- Q
  \gterm{Quantum (plural: quanta)\ {\footnotesize\textcolor{hscmuted}{\mdseries(3)}}}{one packet of energy, of size $E = hf$}{量子}
  % ---- R
  \gterm{Redshift\ {\footnotesize\textcolor{hscmuted}{\mdseries(1, 4)}}}{spectral lines moved to longer wavelengths, for example because the source moves away from us}{红移}
  \gterm{Refractive index ($n$)\ {\footnotesize\textcolor{hscmuted}{\mdseries(2)}}}{how many times slower light is in a material than in a vacuum}{折射率}
  \gterm{Relativistic momentum ($p_\mathrm{v}$)\ {\footnotesize\textcolor{hscmuted}{\mdseries(4)}}}{the momentum that is correct at any speed (on the formula sheet); it grows without limit as $v$ approaches $c$}{相对论动量}
  \gterm{Relativity of simultaneity\ {\footnotesize\textcolor{hscmuted}{\mdseries(4)}}}{events at different places that are simultaneous in one inertial frame are not simultaneous in another}{同时的相对性}
  \gterm{Rest mass ($m_0$)\ {\footnotesize\textcolor{hscmuted}{\mdseries(4)}}}{the mass measured by an observer for whom the object is at rest}{静止质量}
  % ---- S
  \gterm{Saturation current\ {\footnotesize\textcolor{hscmuted}{\mdseries(3)}}}{the largest photocurrent, when every released electron reaches the anode}{饱和电流（饱和光电流）}
  \gterm{Special relativity\ {\footnotesize\textcolor{hscmuted}{\mdseries(4)}}}{Einstein's theory of how space and time behave for observers in inertial frames}{狭义相对论}
  \gterm{Speed of light ($c$)\ {\footnotesize\textcolor{hscmuted}{\mdseries(1)}}}{the speed of all EM waves in a vacuum, exactly \qty{299792458}{m.s^{-1}}}{光速}
  \gterm{Stopping voltage ($V_0$)\ {\footnotesize\textcolor{hscmuted}{\mdseries(3)}}}{the smallest reverse voltage that stops every photoelectron, so the current is zero}{遏止电压}
  \gterm{Superposition\ {\footnotesize\textcolor{hscmuted}{\mdseries(2)}}}{when two waves meet, their displacements add}{叠加（波的叠加）}
  % ---- T
  \gterm{Thought experiment\ {\footnotesize\textcolor{hscmuted}{\mdseries(4)}}}{an imagined experiment, worked through with logic, to see what a theory predicts}{思想实验}
  \gterm{Threshold frequency ($f_0$)\ {\footnotesize\textcolor{hscmuted}{\mdseries(3)}}}{the lowest frequency of light that releases electrons from a metal}{截止频率（极限频率）}
  \gterm{Threshold wavelength ($\lambda_0$)\ {\footnotesize\textcolor{hscmuted}{\mdseries(3)}}}{the longest wavelength of light that releases electrons from a metal}{截止波长（极限波长）}
  \gterm{Time dilation\ {\footnotesize\textcolor{hscmuted}{\mdseries(4)}}}{a moving clock is measured to run slow}{时间膨胀（时间延缓）}
  \gterm{Transmission axis\ {\footnotesize\textcolor{hscmuted}{\mdseries(2)}}}{the direction of vibration that a polarising filter lets through}{透振方向（透光轴）}
  \gterm{Transverse wave\ {\footnotesize\textcolor{hscmuted}{\mdseries(1, 2)}}}{a wave whose vibrations are at \ang{90} to its direction of travel}{横波}
  % ---- U
  \gterm{Ultraviolet catastrophe\ {\footnotesize\textcolor{hscmuted}{\mdseries(3)}}}{the classical prediction of infinite intensity at short wavelengths, which does not happen}{紫外灾难}
  \gterm{Unification\ {\footnotesize\textcolor{hscmuted}{\mdseries(1)}}}{showing that two things that look different are parts of one thing}{统一}
  \gterm{Unpolarised light\ {\footnotesize\textcolor{hscmuted}{\mdseries(2)}}}{light that vibrates in all directions across the beam, changing at random}{非偏振光（自然光）}
  % ---- W
  \gterm{Wavefront\ {\footnotesize\textcolor{hscmuted}{\mdseries(2)}}}{a line joining points of a wave that are in step, such as all the crests}{波阵面（波面）}
  \gterm{Wave--particle duality\ {\footnotesize\textcolor{hscmuted}{\mdseries(3)}}}{light acts like a wave in some experiments and like particles in others}{波粒二象性}
  \gterm{Wien's displacement law\ {\footnotesize\textcolor{hscmuted}{\mdseries(1, 3)}}}{the peak wavelength of a black body is inversely proportional to its absolute temperature: $\lambda_\mathrm{max} = b/T$}{维恩位移定律}
  \gterm{Work function ($\phi$)\ {\footnotesize\textcolor{hscmuted}{\mdseries(3)}}}{the smallest energy needed to remove an electron from a metal's surface}{逸出功}
\end{glossarytable}}
